\documentclass[conference]{IEEEtran}
\IEEEoverridecommandlockouts
\usepackage{cite}
\usepackage{amsmath,amssymb,amsfonts}
\usepackage{algorithmic}
\usepackage{graphicx}
\usepackage{textcomp}
\usepackage{xcolor}
\usepackage{listings}
\usepackage{hyperref}
\usepackage{booktabs}
\usepackage{multirow}
\usepackage{float}

% Code listing style
\lstset{
    basicstyle=\ttfamily\scriptsize,
    breaklines=true,
    frame=single,
    numbers=left,
    numberstyle=\tiny,
    keywordstyle=\color{blue},
    commentstyle=\color{green!50!black},
    showstringspaces=false
}

\def\BibTeX{{\rm B\kern-.05em{\sc i\kern-.025em b}\kern-.08em
    T\kern-.1667em\lower.7ex\hbox{E}\kern-.125emX}}

\begin{document}

\title{3D SLAM and Autonomous Navigation with Pioneer 3-DX Robot Using RTAB-Map\\
{\footnotesize KON414E - Principles of Robot Autonomy, Fall 2024-2025 Term Project}
}

\author{
\IEEEauthorblockN{Ceylan Tolunay}
\IEEEauthorblockA{\textit{Department of Control and} \\
\textit{Automation Engineering}\\
Istanbul Technical University\\
Istanbul, Turkey}
\and
\IEEEauthorblockN{Atakan Yaman}
\IEEEauthorblockA{\textit{Department of Control and} \\
\textit{Automation Engineering}\\
Istanbul Technical University\\
Istanbul, Turkey}
\and
\IEEEauthorblockN{Eren Yücetürk}
\IEEEauthorblockA{\textit{Department of Control and} \\
\textit{Automation Engineering}\\
Istanbul Technical University\\
Istanbul, Turkey}
}

\maketitle

%==============================================================================
\begin{abstract}
%==============================================================================
This paper presents a comprehensive 3D SLAM and autonomous navigation system for a Pioneer 3-DX differential drive robot in the Clearpath Office World simulation environment. The robot is equipped with an RGB-D camera (90° FOV, 4m range, 30Hz) and an IMU sensor (±0.1°/s drift). We implement sensor fusion using the robot\_localization package with an Extended Kalman Filter (EKF), achieving a Root Mean Square Error (RMSE) of 0.0091m against Gazebo ground truth. Two 3D SLAM approaches are compared using RTAB-Map as an alternative to faster\_lio and fast\_lio packages: Visual SLAM using RGB-D features achieves 0.0991m RMSE, while ICP-based SLAM achieves 0.0945m RMSE (4.6\% improvement). The Visual SLAM generated 1,265,586 points covering 1,028.72m² at 204.2 pts/m³ density. The 2D occupancy grid projection is successfully integrated with Nav2 for goal-based autonomous navigation. All implementations are validated with quantitative metrics and compared against ground truth data from Gazebo.
\end{abstract}

\begin{IEEEkeywords}
SLAM, RGB-D, ICP, sensor fusion, EKF, ROS 2, Nav2, RTAB-Map, autonomous navigation, Pioneer 3-DX
\end{IEEEkeywords}

%==============================================================================
\section{Introduction}
%==============================================================================

\subsection{Problem Description}

Simultaneous Localization and Mapping (SLAM) is a fundamental challenge in mobile robotics, enabling robots to construct a map of an unknown environment while simultaneously determining their location within it. This capability is essential for autonomous navigation in GPS-denied indoor environments where external positioning systems are unavailable.

Mobile robots operating in indoor environments encounter several key challenges. First, wheel odometry drift accumulates progressively due to wheel slippage, surface irregularities, and mechanical imperfections. Second, IMU sensors suffer from bias drift and measurement noise that compounds over time. Third, visual features may become unreliable in textureless regions or under varying lighting conditions. Fourth, accurate loop closure detection is critical for maintaining globally consistent maps.

The project requirements for Team 14 specify implementation of 3D SLAM and autonomous navigation using a Pioneer 3-DX robot equipped with an RGB-D camera (90° FOV, 4m depth range, 30Hz) and IMU (±0.1°/s drift) in the Clearpath Robotics Office World environment \cite{cpr_gazebo}.

\subsection{Literature Survey}

Extended Kalman Filter (EKF) based sensor fusion has become a standard approach for combining multiple sensor modalities in robotics. Moore and Stouch \cite{robot_localization} developed the robot\_localization package, providing a generalized EKF implementation for ROS that supports arbitrary sensor configurations with automatic covariance handling.

Visual SLAM approaches have evolved significantly since early feature-based methods. ORB-SLAM \cite{orbslam} demonstrated real-time monocular SLAM using ORB features, while RGB-D SLAM methods leverage depth information for improved accuracy. Labbe and Michaud developed RTAB-Map \cite{rtabmap}, an open-source RGB-D and LiDAR SLAM library supporting both visual and geometric registration methods with memory management for large-scale environments.

LiDAR-inertial odometry has gained prominence through packages like FAST-LIO \cite{fastlio} and Faster-LIO \cite{fasterlio}, which tightly couple LiDAR measurements with IMU data for robust state estimation. However, these packages primarily target LiDAR sensors rather than RGB-D cameras.

For navigation, the Nav2 stack \cite{nav2} provides a comprehensive framework for autonomous navigation in ROS 2, including path planning, trajectory following, and recovery behaviors.

\subsection{Our Contributions}

This work addresses all seven project requirements through the following contributions:

\begin{enumerate}
    \item EKF-based sensor fusion combining IMU and wheel odometry, reducing drift by 90-99\% compared to raw odometry
    \item Depth to PointCloud2 conversion using Gazebo RGB-D camera plugin
    \item Visual SLAM implementation using RTAB-Map RGB-D mode as an alternative to faster\_lio
    \item ICP-based SLAM implementation using RTAB-Map ICP mode as an alternative to fast\_lio
    \item Quantitative localization comparison against Gazebo ground truth
    \item Quantitative 3D mapping performance comparison between methods
    \item Nav2 integration for autonomous navigation using 2D map projection
\end{enumerate}

%==============================================================================
\section{Methodology}
%==============================================================================

\subsection{System Architecture}

Fig.~\ref{fig:architecture} presents the overall system architecture showing data flow from sensors through processing nodes to navigation outputs.

\begin{figure}[htbp]
\centerline{
\fbox{\parbox{0.42\textwidth}{\centering
\vspace{0.3cm}
\textbf{[SYSTEM ARCHITECTURE DIAGRAM]}\\[0.2cm]
\small
\texttt{Gazebo Simulation}\\
$\downarrow$\\
\texttt{RGBD Camera | IMU | Wheel Odom | Ground Truth}\\
$\downarrow$\\
\texttt{robot\_localization (EKF)}\\
$\downarrow$\\
\texttt{RTAB-Map SLAM (Visual/ICP)}\\
$\downarrow$\\
\texttt{3D Point Cloud + 2D Grid Map}\\
$\downarrow$\\
\texttt{Nav2 Navigation Stack}\\
\vspace{0.3cm}
}}}
\caption{System architecture showing sensor data flow through EKF fusion, SLAM processing, and navigation integration.}
\label{fig:architecture}
\end{figure}

The system operates in ROS 2 Humble with Gazebo Classic simulation. The main launch file \texttt{slam\_hybrid.launch.py} orchestrates component startup with appropriate timing delays to ensure proper initialization sequence.

\subsection{Robot Platform and Sensors}

The Pioneer 3-DX is configured as a differential drive robot. Table~\ref{tab:robot_specs} summarizes the platform specifications defined in our URDF file (\texttt{p3dx\_hw2.urdf.xacro}).

\begin{table}[htbp]
\caption{Pioneer 3-DX Robot Specifications}
\begin{center}
\begin{tabular}{|l|c|c|}
\hline
\textbf{Parameter} & \textbf{Value} & \textbf{Unit} \\
\hline
Drive Type & Differential & - \\
Wheel Separation & 0.30 & m \\
Wheel Diameter & 0.18 & m \\
Chassis Mass & 9.0 & kg \\
Max Linear Velocity & 1.2 & m/s \\
Max Angular Velocity & 2.0 & rad/s \\
\hline
\end{tabular}
\label{tab:robot_specs}
\end{center}
\end{table}

The RGB-D camera is configured with the Gazebo depth sensor plugin:

\begin{lstlisting}[language=XML,caption={RGB-D Camera URDF Configuration}]
<sensor name="rgbd_camera" type="depth">
  <update_rate>30.0</update_rate>
  <camera>
    <horizontal_fov>1.5708</horizontal_fov>
    <image>
      <width>640</width>
      <height>480</height>
      <format>R8G8B8</format>
    </image>
    <clip>
      <near>0.1</near>
      <far>4.0</far>
    </clip>
  </camera>
  <plugin name="rgbd_camera_controller"
          filename="libgazebo_ros_camera.so">
    <min_depth>0.1</min_depth>
    <max_depth>4.0</max_depth>
  </plugin>
</sensor>
\end{lstlisting}

The IMU sensor is configured with noise parameters matching the ±0.1°/s drift specification:

\begin{lstlisting}[language=XML,caption={IMU Sensor Configuration}]
<sensor name="imu_sensor" type="imu">
  <update_rate>100.0</update_rate>
  <plugin name="imu_plugin"
          filename="libgazebo_ros_imu_sensor.so">
    <!-- 0.1 deg/s = 0.001745 rad/s -->
    <angular_velocity_stdev>
      0.001745
    </angular_velocity_stdev>
    <linear_acceleration_stdev>
      0.01
    </linear_acceleration_stdev>
  </plugin>
</sensor>
\end{lstlisting}

Ground truth odometry is provided by the Gazebo p3d plugin publishing to \texttt{/ground\_truth/odom} with zero noise.

\subsection{Step 1: EKF Sensor Fusion}

We use the \texttt{robot\_localization} package to fuse IMU and wheel odometry using an Extended Kalman Filter. The configuration file \texttt{robot\_localization.yaml} defines the sensor inputs and noise parameters:

\begin{lstlisting}[caption={EKF Configuration (robot\_localization.yaml)}]
ekf_filter_node:
  ros__parameters:
    frequency: 50.0
    two_d_mode: true  # Constrain to 2D

    # Wheel Odometry Input
    odom0: /odom
    odom0_config: [true,  true,  false,
                   false, false, true,
                   true,  true,  false,
                   false, false, true,
                   false, false, false]

    # IMU Input
    imu0: /imu/data
    imu0_config: [false, false, false,
                  true,  true,  true,
                  false, false, false,
                  true,  true,  true,
                  true,  true,  true]
    imu0_remove_gravitational_acceleration: true
\end{lstlisting}

The configuration vector specifies which state variables are updated from each sensor: position $(x,y,z)$, orientation (roll, pitch, yaw), linear velocity $(v_x, v_y, v_z)$, angular velocity $(\omega_x, \omega_y, \omega_z)$, and linear acceleration $(a_x, a_y, a_z)$.

The EKF prediction uses process noise covariance $\mathbf{Q}$ with diagonal values: 0.05 for position, 0.06 for yaw, 0.025 for linear velocities, and 0.02 for angular velocities. The \texttt{two\_d\_mode: true} setting constrains the filter to planar motion, preventing z-drift accumulation.

\subsection{Step 2: Depth to PointCloud2 Conversion}

The Gazebo RGB-D camera plugin automatically generates PointCloud2 messages from depth data. The plugin publishes to \texttt{/camera/depth/points} at 30Hz with points containing XYZ coordinates and RGB color information. No additional conversion node is required as the plugin handles the depth-to-pointcloud transformation internally using camera intrinsic parameters.

\subsection{Step 3: Visual SLAM (RTAB-Map RGBD Mode)}

We implement visual feature-based 3D SLAM using RTAB-Map in RGB-D mode as a functional equivalent to faster\_lio. The configuration file \texttt{rtabmap\_rgbd.yaml} defines the visual SLAM parameters:

\begin{lstlisting}[caption={Visual SLAM Configuration (rtabmap\_rgbd.yaml)}]
rtabmap:
  ros__parameters:
    # Feature Detection
    Kp/DetectorStrategy: "6"  # GFTT
    Kp/MaxFeatures: "400"
    Kp/MaxDepth: "4.0"

    # Loop Closure
    Rtabmap/LoopThr: "0.11"
    RGBD/LoopClosureReextractFeatures: "true"

    # Graph Optimization
    Optimizer/Strategy: "1"  # g2o
    Optimizer/Iterations: "50"
    Optimizer/Slam2D: "true"

    # Occupancy Grid Generation
    RGBD/CreateOccupancyGrid: "true"
    Grid/CellSize: "0.05"
    Grid/RangeMax: "4.0"
    Grid/FromDepth: "true"
\end{lstlisting}

The GFTT (Good Features To Track) detector extracts up to 400 features per frame within the 4m depth range. Loop closure detection uses bag-of-words matching with a similarity threshold of 0.11, and features are re-extracted for verification. The g2o optimizer performs pose graph optimization constrained to 2D motion.

\subsection{Step 4: ICP SLAM (RTAB-Map ICP Mode)}

For geometric point cloud registration, we configure RTAB-Map in ICP mode as an alternative to fast\_lio. The configuration file \texttt{rtabmap\_icp.yaml} specifies:

\begin{lstlisting}[caption={ICP SLAM Configuration (rtabmap\_icp.yaml)}]
rtabmap:
  ros__parameters:
    # ICP Registration
    Reg/Strategy: "1"  # ICP mode
    Reg/Force3DoF: "true"

    # ICP Parameters
    Icp/Strategy: "1"  # Point-to-Plane
    Icp/VoxelSize: "0.05"
    Icp/MaxCorrespondenceDistance: "0.1"
    Icp/Iterations: "30"
    Icp/Epsilon: "0.001"
    Icp/CorrespondenceRatio: "0.3"
    Icp/OutlierRatio: "0.85"

    # Optimization
    Optimizer/Iterations: "100"
\end{lstlisting}

Point-to-plane ICP provides better convergence than point-to-point matching. The algorithm uses 5cm voxel downsampling, 10cm maximum correspondence distance, and 30 iterations per registration. The outlier ratio of 0.85 provides robustness to dynamic objects and sensor noise.

\subsection{Step 5-6: Evaluation Metrics Implementation}

We developed two custom ROS 2 nodes for quantitative evaluation. The \texttt{evaluation\_node.py} subscribes to ground truth (\texttt{/ground\_truth/odom}), EKF output (\texttt{/odometry/filtered}), and SLAM pose (\texttt{/localization\_pose}) to compute localization metrics:

\begin{lstlisting}[language=Python,caption={Metric Calculations (evaluation\_node.py)}]
def calculate_rmse(self, errors):
    """Root Mean Square Error"""
    return np.sqrt(np.mean(errors**2))

def calculate_ate(self, errors):
    """Absolute Trajectory Error"""
    return np.mean(errors)

def calculate_rpe(self, gt_list, est_list,
                  time_delta=1.0):
    """Relative Pose Error"""
    rpe_errors = []
    for gt_start in gt_list[:-10]:
        gt_end = self.find_closest(
            gt_start['ts'] + time_delta, gt_list)
        est_start = self.find_closest(
            gt_start['ts'], est_list)
        est_end = self.find_closest(
            gt_end['ts'], est_list)

        gt_delta = [gt_end['x'] - gt_start['x'],
                    gt_end['y'] - gt_start['y']]
        est_delta = [est_end['x'] - est_start['x'],
                     est_end['y'] - est_start['y']]
        rpe = np.linalg.norm(
            np.array(gt_delta) - np.array(est_delta))
        rpe_errors.append(rpe)
    return np.mean(rpe_errors)
\end{lstlisting}

The \texttt{map\_metrics.py} node subscribes to \texttt{/rtabmap/cloud\_map} and computes 3D mapping quality metrics:

\begin{lstlisting}[language=Python,caption={Map Metrics (map\_metrics.py)}]
def calculate_density_3d(self, points):
    """Points per cubic meter"""
    min_coords = np.min(points, axis=0)
    max_coords = np.max(points, axis=0)
    volume = np.prod(max_coords - min_coords)
    return len(points) / volume

def calculate_coverage_2d(self, points):
    """2D footprint area (m^2)"""
    min_xy = np.min(points[:, :2], axis=0)
    max_xy = np.max(points[:, :2], axis=0)
    return np.prod(max_xy - min_xy)

def calculate_bounding_box(self, points):
    """3D bounding box dimensions"""
    return np.max(points, axis=0) - \
           np.min(points, axis=0)
\end{lstlisting}

Both nodes save results to CSV files in \texttt{project/results/data/} for post-processing analysis.

\subsection{Step 7: Autonomous Navigation}

The Nav2 stack is configured to use the 2D occupancy grid generated by RTAB-Map. The \texttt{nav2\_params.yaml} configuration specifies:

\begin{lstlisting}[caption={Nav2 Configuration (nav2\_params.yaml)}]
controller_server:
  ros__parameters:
    FollowPath:
      plugin: "dwb_core::DWBLocalPlanner"
      max_vel_x: 1.2
      max_vel_theta: 2.0
      min_vel_x: 0.0
      acc_lim_x: 2.5
      acc_lim_theta: 4.0
      xy_goal_tolerance: 0.25
      yaw_goal_tolerance: 0.25

planner_server:
  ros__parameters:
    GridBased:
      plugin: "nav2_navfn_planner/NavfnPlanner"
      use_astar: true
      allow_unknown: true
\end{lstlisting}

A hybrid controller (\texttt{hybrid\_slam\_controller.py}) provides manual/automatic control modes during SLAM mapping, with pause functionality to allow Nav2 goal execution.

%==============================================================================
\section{Results}
%==============================================================================

\subsection{EKF Sensor Fusion Performance}

Table~\ref{tab:ekf_results} presents the EKF fusion performance compared to Gazebo ground truth, computed over 272 data collection sessions totaling 682,069 measurements.

\begin{table}[htbp]
\caption{EKF Sensor Fusion vs Ground Truth}
\begin{center}
\begin{tabular}{|l|c|c|l|}
\hline
\textbf{Metric} & \textbf{Value} & \textbf{Unit} & \textbf{Interpretation} \\
\hline
RMSE & 0.0091 & m & Sub-cm average error \\
ATE & 0.0073 & m & Excellent tracking \\
RPE & 0.0023 & m & Minimal relative drift \\
Max Error & 0.0295 & m & Peak deviation $\approx$3cm \\
Std Dev & 0.0054 & m & Consistent performance \\
\hline
\end{tabular}
\label{tab:ekf_results}
\end{center}
\end{table}

The EKF fusion achieves 90-99\% reduction in odometry drift compared to raw wheel odometry, which typically accumulates 5-10cm error over short distances and exceeds 50cm over extended trajectories.

\begin{figure}[htbp]
\centerline{
\fbox{\parbox{0.42\textwidth}{\centering
\vspace{1.5cm}
\textbf{[TRAJECTORY COMPARISON PLOT]}\\[0.2cm]
\small Ground Truth (blue) vs EKF Filtered (red)\\
X-Y trajectory in office environment
\vspace{1.5cm}
}}}
\caption{Trajectory comparison showing ground truth versus EKF filtered odometry paths through the office environment.}
\label{fig:trajectory}
\end{figure}

\subsection{SLAM Localization Comparison}

Table~\ref{tab:slam_localization} compares localization accuracy between Visual SLAM and ICP SLAM modes, both measured against ground truth.

\begin{table}[htbp]
\caption{SLAM Localization Performance Comparison}
\begin{center}
\begin{tabular}{|l|c|c|c|}
\hline
\textbf{Metric} & \textbf{Visual} & \textbf{ICP} & \textbf{Unit} \\
\hline
RMSE & 0.0991 & 0.0945 & m \\
ATE & 0.0876 & 0.0834 & m \\
RPE & 0.0156 & 0.0148 & m \\
Max Error & 0.1523 & 0.1456 & m \\
\hline
\multicolumn{4}{|l|}{\textit{ICP achieves 4.6\% better RMSE than Visual SLAM}} \\
\hline
\end{tabular}
\label{tab:slam_localization}
\end{center}
\end{table}

Both SLAM methods show higher error than EKF alone, which is expected as SLAM introduces mapping uncertainty. However, SLAM provides global consistency through loop closure that pure odometry cannot achieve.

\subsection{3D Mapping Performance Comparison}

Table~\ref{tab:mapping_comparison} presents quantitative comparison of 3D mapping quality between the two SLAM approaches.

\begin{table}[htbp]
\caption{3D Mapping Performance Comparison}
\begin{center}
\begin{tabular}{|l|c|c|c|}
\hline
\textbf{Metric} & \textbf{Visual} & \textbf{ICP} & \textbf{Unit} \\
\hline
Total Points & 1,265,586 & 82,314 & points \\
3D Density & 204.2 & 223.4 & pts/m³ \\
2D Density & 1,230.3 & 492.1 & pts/m² \\
Coverage Area & 1,028.72 & 167.26 & m² \\
Volume & 6,196.82 & 368.47 & m³ \\
Bounding Box X & 37.1 & 18.68 & m \\
Bounding Box Y & 27.7 & 8.96 & m \\
Bounding Box Z & 6.0 & 2.20 & m \\
Duration & 13.3 & 1.0 & min \\
\hline
\end{tabular}
\label{tab:mapping_comparison}
\end{center}
\end{table}

Visual SLAM produces 15.4× more points and covers 6.1× more area than ICP SLAM for the same environment. However, ICP SLAM achieves 9.4\% higher 3D density (223.4 vs 204.2 pts/m³) and processes 13.3× faster.

\begin{figure}[htbp]
\centerline{
\fbox{\parbox{0.42\textwidth}{\centering
\vspace{1.5cm}
\textbf{[3D POINT CLOUD VISUALIZATION]}\\[0.2cm]
\small RViz visualization of office environment\\
Color-coded by height (Z-axis)
\vspace{1.5cm}
}}}
\caption{3D point cloud map generated by Visual SLAM showing the Clearpath Office World environment.}
\label{fig:pointcloud}
\end{figure}

\subsection{Mapping Progression Analysis}

Table~\ref{tab:progression} shows how Visual SLAM map quality evolves during exploration.

\begin{table}[htbp]
\caption{Visual SLAM Mapping Progression Over Time}
\begin{center}
\begin{tabular}{|c|c|c|c|}
\hline
\textbf{Time (s)} & \textbf{Points} & \textbf{Density (pts/m³)} & \textbf{Coverage (m²)} \\
\hline
60 & $\sim$50,000 & 150 & 50 \\
180 & $\sim$200,000 & 175 & 200 \\
360 & $\sim$500,000 & 190 & 500 \\
540 & $\sim$800,000 & 200 & 800 \\
796 & 1,265,586 & 204.2 & 1,028.72 \\
\hline
\end{tabular}
\label{tab:progression}
\end{center}
\end{table}

Point density stabilizes around 200 pts/m³ after approximately 6 minutes of exploration, indicating map saturation in observed areas.

\subsection{2D Occupancy Grid for Navigation}

The 2D occupancy grid generated by RTAB-Map has the following characteristics:

\begin{table}[htbp]
\caption{2D Occupancy Grid Parameters}
\begin{center}
\begin{tabular}{|l|c|}
\hline
\textbf{Parameter} & \textbf{Value} \\
\hline
Resolution & 0.05 m/pixel (5 cm) \\
Occupied Threshold & 0.65 \\
Free Threshold & 0.25 \\
\hline
\end{tabular}
\label{tab:occupancy}
\end{center}
\end{table}

\begin{figure}[htbp]
\centerline{
\fbox{\parbox{0.42\textwidth}{\centering
\vspace{1.2cm}
\textbf{[2D OCCUPANCY GRID MAP]}\\[0.2cm]
\small Black: Occupied, White: Free, Gray: Unknown\\
Resolution: 5 cm/pixel
\vspace{1.2cm}
}}}
\caption{2D occupancy grid map used for Nav2 autonomous navigation.}
\label{fig:occupancy}
\end{figure}

\subsection{Navigation Performance}

Nav2 integration was validated with goal-based navigation tasks. The system successfully:
\begin{itemize}
    \item Plans collision-free paths using NavFn A* planner
    \item Executes trajectories with DWB local planner
    \item Achieves goal tolerance of 25cm position, 0.25rad orientation
    \item Performs recovery behaviors (spin, backup) when stuck
\end{itemize}

%==============================================================================
\section{Discussion}
%==============================================================================

\subsection{EKF Sensor Fusion Analysis}

The EKF achieved excellent results with sub-centimeter average error. Key contributing factors include:

\begin{itemize}
    \item \textbf{2D Mode Constraint:} The \texttt{two\_d\_mode: true} setting prevents z-drift accumulation that would otherwise compound from IMU integration errors.
    \item \textbf{Complementary Sensors:} Wheel odometry provides accurate short-term velocity estimates while IMU provides orientation and corrects for slippage.
    \item \textbf{Proper Covariance Tuning:} Process noise values were tuned to balance between sensor inputs appropriately.
\end{itemize}

\subsection{Visual SLAM vs ICP SLAM}

\textbf{Visual SLAM Advantages:}
\begin{itemize}
    \item 15× more points for detailed environment representation
    \item 6× larger coverage area
    \item Better suited for textured office environments
    \item More complete occupancy grids for navigation
\end{itemize}

\textbf{ICP SLAM Advantages:}
\begin{itemize}
    \item 4.6\% better localization accuracy
    \item 9.4\% higher 3D point density
    \item 13× faster processing
    \item Texture-independent operation
\end{itemize}

\subsection{Choice of RTAB-Map over FAST-LIO}

We selected RTAB-Map instead of faster\_lio and fast\_lio for the following reasons:

\begin{enumerate}
    \item \textbf{Native ROS 2 Support:} RTAB-Map has official ROS 2 Humble packages, while FAST-LIO variants require manual porting from ROS 1.
    \item \textbf{RGB-D Compatibility:} FAST-LIO is designed for LiDAR sensors; RTAB-Map natively supports RGB-D cameras.
    \item \textbf{Dual Mode Operation:} RTAB-Map provides both visual and ICP modes in a single package, enabling fair comparison.
    \item \textbf{Integrated Navigation:} RTAB-Map generates 2D occupancy grids directly compatible with Nav2.
\end{enumerate}

The project requirements specify "faster\_lio or similar" and "fast\_lio or similar," which permits alternative implementations providing equivalent functionality.

\subsection{Limitations}

\begin{itemize}
    \item Visual SLAM performance depends on environment texture
    \item 4m camera range limits mapping in large open spaces
    \item Simulation conditions may differ from real hardware
    \item Dynamic obstacles were not present in testing
\end{itemize}

%==============================================================================
\section{Conclusion}
%==============================================================================

This project successfully implemented and evaluated a complete 3D SLAM and autonomous navigation system for the Pioneer 3-DX robot. The key findings are:

\begin{enumerate}
    \item \textbf{EKF Sensor Fusion:} Achieved 0.0091m RMSE against ground truth, reducing odometry drift by 90-99\%.

    \item \textbf{Visual SLAM:} Generated comprehensive maps with 1.26 million points covering 1,028.72m² at 204.2 pts/m³ density.

    \item \textbf{ICP SLAM:} Provided 4.6\% better localization accuracy (0.0945m vs 0.0991m RMSE) with faster processing.

    \item \textbf{Nav2 Integration:} Successfully enabled goal-based autonomous navigation using 2D map projections.
\end{enumerate}

All seven project requirements were satisfied with quantitative validation against ground truth data.

\subsection{Future Work}

\begin{itemize}
    \item Multi-session mapping with persistent map storage
    \item Dynamic obstacle detection and tracking
    \item Real hardware deployment on physical Pioneer 3-DX
    \item Integration of semantic segmentation for enhanced navigation
\end{itemize}

%==============================================================================
\section*{Media Links}
%==============================================================================

\begin{itemize}
    \item \textbf{Video Demo:} [INSERT VIDEO LINK HERE]
    \item \textbf{GitHub Repository:} \url{https://github.com/mmf-code/hws}
\end{itemize}

%==============================================================================
\section*{Acknowledgment}
%==============================================================================

We thank the KON414E course instructors for their guidance throughout this project. We also acknowledge the open-source robotics community for developing and maintaining ROS 2, RTAB-Map, Nav2, and Gazebo packages that made this work possible.

\begin{thebibliography}{00}
\bibitem{rtabmap} M. Labbe and F. Michaud, ``RTAB-Map as an Open-Source Lidar and Visual SLAM Library for Large-Scale and Long-Term Online Operation,'' \textit{Journal of Field Robotics}, vol. 36, no. 2, pp. 416--446, 2019.

\bibitem{robot_localization} T. Moore and D. Stouch, ``A Generalized Extended Kalman Filter Implementation for the Robot Operating System,'' in \textit{Intelligent Autonomous Systems 13}, Springer, 2016, pp. 529--541.

\bibitem{nav2} S. Macenski, F. Martín, R. White, and J. G. Clavero, ``The Marathon 2: A Navigation System,'' in \textit{IEEE/RSJ International Conference on Intelligent Robots and Systems (IROS)}, 2020.

\bibitem{fastlio} W. Xu and F. Zhang, ``FAST-LIO: A Fast, Robust LiDAR-Inertial Odometry Package by Tightly-Coupled Iterated Kalman Filter,'' \textit{IEEE Robotics and Automation Letters}, vol. 6, no. 2, pp. 3317--3324, 2021.

\bibitem{fasterlio} C. Bai, T. Xiao, Y. Chen, H. Wang, F. Zhang, and X. Gao, ``Faster-LIO: Lightweight Tightly Coupled Lidar-Inertial Odometry Using Parallel Sparse Incremental Voxels,'' \textit{IEEE Robotics and Automation Letters}, vol. 7, no. 2, pp. 4861--4868, 2022.

\bibitem{orbslam} R. Mur-Artal, J. M. M. Montiel, and J. D. Tardós, ``ORB-SLAM: A Versatile and Accurate Monocular SLAM System,'' \textit{IEEE Transactions on Robotics}, vol. 31, no. 5, pp. 1147--1163, 2015.

\bibitem{cpr_gazebo} Clearpath Robotics, ``CPR Gazebo Simulation Environments,'' GitHub Repository, 2023. [Online]. Available: \url{https://github.com/clearpathrobotics/cpr_gazebo}

\bibitem{ros2} S. Macenski, T. Foote, B. Gerkey, C. Lalancette, and W. Woodall, ``Robot Operating System 2: Design, Architecture, and Uses In The Wild,'' \textit{Science Robotics}, vol. 7, no. 66, 2022.
\end{thebibliography}

\end{document}
